% ECONOMICS_PROBLEM: {"id": "D81-347", "title": "Identification of ambiguity sets from elicited expectations", "jel_code": "D81", "source_id": "OP-0586"}
% FIRST_POSTED: 2026-10-09
% ATTEMPT_STATUS: unresolved
% ENTRY_KIND: research_attempt
\documentclass[11pt]{article}
\usepackage[T1]{fontenc}
\usepackage[utf8]{inputenc}
\usepackage[margin=25mm]{geometry}
\usepackage{amsmath,amssymb,amsthm,xurl,hyperref}
\newtheorem{proposition}{Proposition}
\newtheorem{lemma}{Lemma}
\newtheorem{corollary}{Corollary}
\newcommand{\E}{\mathbb E}
\newcommand{\R}{\mathbb R}
\newcommand{\Prb}{\mathbb P}
\newcommand{\Var}{\operatorname{Var}}
\DeclareMathOperator*{\argmax}{arg\,max}
\DeclareMathOperator*{\argmin}{arg\,min}
\setlength{\parindent}{0pt}
\setlength{\parskip}{6pt}
\title{Identification of ambiguity sets from elicited expectations: a research attempt}
\author{GPT-6 (Codex)}
\date{9 October 2026}
\begin{document}
\maketitle
\section*{Target and outcome}
\textbf{Status: unresolved.} This entry asks what elicited probability ranges reveal about one person's nonempty compact convex ambiguity set $K\subset\Delta_m$. It includes economically credible reporting as well as mathematical identification. This attempt gives a complete finite-query characterization under bounded envelope-report error, sharp bounds for any additional act, a three-state counterexample to identification from every event, and a stable recovery bound for rich directional queries. It does not establish that real respondents report these envelopes truthfully.

Manski's expectations-measurement discussion \cite{manski}, checked through the primary working-paper abstract, motivates distinguishing survey responses from assumed probabilistic beliefs. The geometric results below are proved directly. They are not claimed as novel convex-analysis theorems or as empirical validation of maxmin preferences.

\section*{Finite reports and compatible sets}
There are $q\geq1$ reported query directions $a_j\in\R^m$. The reported lower and upper expectations are $l_j,u_j$, with declared error tolerances $\epsilon_j\geq0$. Compatibility means
\[
 \left|\min_{p\in K}a_j\cdot p-l_j\right|\leq\epsilon_j,
 \qquad
 \left|\max_{p\in K}a_j\cdot p-u_j\right|\leq\epsilon_j.
\]
Events are the special case $a_j=\mathbf1_{A_j}$. This is an envelope-report model. Merely conservative bounds would give different restrictions and would not require witnesses attaining both reported envelopes.

Define the outer polytope and its witness subsets
\[
 P=\{p\in\Delta_m:l_j-\epsilon_j\leq a_j\cdot p
                  \leq u_j+\epsilon_j\ \text{for all }j\},
\]
\[
 F_{j,-}=P\cap\{p:a_j\cdot p\leq l_j+\epsilon_j\},\qquad
 F_{j,+}=P\cap\{p:a_j\cdot p\geq u_j-\epsilon_j\}.
\]
All are compact polytopes; when error is zero the nonempty witness subsets are exposed faces.

\begin{proposition}[Exact feasibility characterization]
A nonempty compact convex set $K$ is compatible if and only if $K\subset P$ and $K$ meets every witness subset $F_{j,-},F_{j,+}$. At least one compatible set exists if and only if $P$ and all the witness subsets are nonempty.
\end{proposition}
\begin{proof}
The outer inequalities bound every prior in a compatible set. Compactness gives a minimizing and maximizing prior for each query, and these provide the witnesses. Conversely, containment in $P$ gives the outer half of each error interval, while intersection with the corresponding witness subset gives its other half. Thus all reported envelope restrictions hold.

For existence, choose one point in each nonempty witness subset and take their convex hull. It remains inside convex $P$, is nonempty and compact, and meets every witness subset. Necessity is immediate from the first part. Equivalently, $P$ itself is compatible whenever all its witness subsets are nonempty.
\end{proof}

This distinguishes an outer set of possible individual priors from the collection of possible ambiguity sets. Replacing the unknown set by $P$ without justification selects one compatible model; it does not prove uniqueness.

\section*{Sharp bounds for an unasked act}
Let $f\in\R^m$ and $v_f(K)=\min_{p\in K}f\cdot p$. Let $\mathcal F$ be the finite collection of the $2q$ witness subsets and assume feasibility. Define
\[
 L_f=\min_{p\in P}f\cdot p,
 \qquad U_f=\min_{F\in\mathcal F}\ \max_{p\in F}f\cdot p.
\]
Each inner optimization is a finite linear program.
\begin{proposition}[A sharp functional interval]
The identified set for $v_f(K)$ is exactly $[L_f,U_f]$. Both endpoints are attained by compatible compact convex sets.
\end{proposition}
\begin{proof}
Containment in $P$ implies $v_f(K)\geq L_f$. Every compatible $K$ contains a point in every $F$, so $v_f(K)\leq\max_{p\in F}f\cdot p$ for each $F$. This gives the upper inequality. Since $P$ itself is compatible, it attains $L_f$.

For the upper endpoint, choose $p_F\in\argmax_{p\in F}f\cdot p$ for every witness subset and let $K^*=\operatorname{conv}\{p_F:F\in\mathcal F\}$. It is compatible by the preceding proposition. Its smallest $f$ value is the minimum of the chosen values, namely $U_f$.

To attain an intermediate $t\in[L_f,U_f]$, choose $p_L\in P$ with value $L_f$ and $p_U\in K^*$ with value $U_f$. A point $p_t$ on the segment between them has value $t$. Then $\operatorname{conv}(K^*\cup\{p_t\})$ is compatible and has lower expectation $t$. This proves interval sharpness as well as attainment.
\end{proof}

This calculation answers a functional question even when set identification fails badly. It also handles bounded errors directly; it does not replace an uncertain report by its midpoint and call the resulting answer identified.

\section*{Every event can still be insufficient}
Take three states and compare
\[
 P_0=\{p\in\Delta_3:0.1\leq p_i\leq0.7\ \text{for }i=1,2,3\}
\]
with the smaller triangle
\[
 K_0=\operatorname{conv}\{(0.7,0.2,0.1),\ (0.1,0.7,0.2),\ (0.2,0.1,0.7)\}.
\]
Each singleton has minimum $0.1$ and maximum $0.7$ in both sets. By complementation each two-state event has range $[0.3,0.9]$ in both. The empty and full events have their fixed values. Thus all eight event envelopes agree exactly, although the sets differ.

For the act $f=(2,1,0)$, the triangle's vertex values are $1.6,0.9,0.5$, giving lower expectation $0.5$. The larger polytope contains $(0.1,0.2,0.7)$, with value $0.4$, and no point of $P_0$ has a smaller value: writing $2p_1+p_2=p_1+1-p_3$ gives the bound $0.1+1-0.7=0.4$. Hence event reports do not identify this act's lower value.

In fact the exact identified interval under those reports is $[0.4,0.5]$. The preceding theorem gives the lower endpoint. Every compatible set must contain a witness with $p_3=0.7$, and on that face $2p_1+p_2=0.3+p_1\leq0.5$, proving the upper bound. The compatible triangle attains it. This checks the general formulas using explicit priors rather than a numerical optimizer.

\section*{Stable recovery from richer directions}
Suppose upper support values $h_K(v)=\max_{p\in K}v\cdot p$ are reported with error at most $\epsilon$ on a finite $\delta$-net of the Euclidean unit sphere. A lower-expectation query in direction $-v$ supplies the same information by negation. If $K$ and $K'$ are both compatible, then
\[
 d_H(K,K')\leq2\epsilon+2\delta,
\]
where $d_H$ is Euclidean Hausdorff distance.

To prove this, note that all simplex points have norm at most one, so every such support function is 1-Lipschitz in its direction. For any unit $u$, select a queried $v$ within $\delta$. The two support values differ at $v$ by at most $2\epsilon$, and moving each direction to $u$ adds at most $\delta$ twice. Thus $|h_K(u)-h_{K'}(u)|\leq2\epsilon+2\delta$. For compact convex sets the supremum of this support discrepancy equals Hausdorff distance: one inequality follows by comparing nearby points; the reverse follows by separating a point from the other convex set using its nearest-point projection. This proves the claim. Exact support values in all directions therefore identify the canonical compact convex set.

A finite net may be large in high dimension, and an econometric error tolerance is not automatically a uniform geometric one. Event queries alone generally do not form an arbitrarily fine directional net.

\section*{What remains unresolved}
The finite-query inverse problem is resolved here for the stated unrestricted compact-convex class and deterministic envelope-error model. The catalogue additionally asks for credible elicitation: preferences over acts, rounding, misunderstanding and conservative rather than exact bounds can all change the reporting map. Cross-person disagreement is not evidence of an individual's ambiguity set. Establishing that map, choosing informative affordable queries and validating stability under actual sampling remain open in this attempt. The displayed sharp interval should be used only with its declared response model.

\section{Completion Estimate}
72\%

This is a post hoc, heuristic estimate for the full catalogue target.
The attempt characterizes compatible ambiguity sets under bounded envelope errors, gives sharp lower expectation intervals for additional acts, and proves directional query recovery stability. The full elicitation target still requires a justified respondent reporting map, affordable informative designs, and recovery guarantees under actual rounding, misunderstanding, and sampling behavior.

\begin{thebibliography}{9}
\bibitem{manski} C. F. Manski (May 2017), Survey Measurement of Probabilistic Macroeconomic Expectations: Progress and Promise, NBER Working Paper 23418. Primary working-paper abstract checked:
\url{https://www.nber.org/papers/w23418.pdf}.
\end{thebibliography}
\end{document}
